\documentclass[12pt,a4paper]{article}

\usepackage{amsmath,amssymb,amsthm}
\usepackage{mathtools}
\usepackage{enumitem}
\usepackage[margin=1in]{geometry}
\usepackage{hyperref}
\usepackage{cleveref}

% Theorem environments
\newtheorem{theorem}{Theorem}[section]
\newtheorem{proposition}[theorem]{Proposition}
\newtheorem{lemma}[theorem]{Lemma}
\newtheorem{corollary}[theorem]{Corollary}
\newtheorem{conjecture}[theorem]{Conjecture}
\theoremstyle{definition}
\newtheorem{definition}[theorem]{Definition}
\newtheorem{example}[theorem]{Example}
\newtheorem{openquestion}[theorem]{Open Question}
\theoremstyle{remark}
\newtheorem{remark}[theorem]{Remark}

% Shortcuts
\newcommand{\Z}{\mathbb{Z}}
\newcommand{\R}{\mathbb{R}}
\newcommand{\C}{\mathbb{C}}
\newcommand{\HH}{\mathbb{H}}
\newcommand{\OO}{\mathbb{O}}
\newcommand{\Sed}{\mathbb{S}}
\newcommand{\Tri}{\mathbb{T}}
\newcommand{\ZSet}{\Z_2\text{-}\mathbf{Set}}
\newcommand{\CD}{\mathrm{CD}}
\DeclareMathOperator{\Aut}{Aut}
\DeclareMathOperator{\Gal}{Gal}
\DeclareMathOperator{\id}{id}
\DeclareMathOperator{\Im}{Im}
\newcommand{\swap}{\mathtt{swap}}

\title{Tower Generation from M\"{o}bius Holonomy:\\
Canonical Resolution of Representational Obstruction}

\author{Adam Kevin Morgan}

\date{\today \\ {\small Draft --- working paper}}

\begin{document}
\maketitle

\begin{abstract}
We show that the M\"{o}bius connection $\omega = \tfrac{1}{2}d\theta$
on $S^1$ generates an infinite canonical tower of algebraic and
geometric structures through iterated resolution of representational
obstruction. At each level, the subobject classifier of the
associated topos contains \emph{swap} values---propositions that are
locally definite but globally unresolvable. These constitute
cohomological obstructions whose minimal isometric resolution is
unique and forced, producing the Cayley--Dickson doubling of the
current algebra.

The first four levels recover the normed division algebras
$\R, \C, \HH, \OO$. Beyond $\OO$, zero divisors appear and the
resolution outputs constraint geometry rather than division algebras.
We prove three structural theorems for the post-division-algebra
regime: (1)~a \emph{two-sidedness theorem} establishing that zero
divisor relations among pure imaginary elements are symmetric, derived
from the linearized norm identity and flexibility of Cayley--Dickson
algebras; (2)~a \emph{degree-doubling theorem} showing that every
zero-divisor node exactly doubles its connectivity at each level of
the tower; and (3)~a \emph{non-termination theorem} establishing
that the tower generates strictly new structure at every level.
Together, these establish that the tower is the unique canonical
object generated by iterating the M\"{o}bius connection's holonomy
through its own representational obstructions, with no choice at
any step.

Explicit computation through six post-NDA levels (dimension~16
to~512, over eight million zero-divisor pairs) reveals
unexpectedly rigid structure: the strata count follows
$2^n - 1$ exactly, singular values of multiplication maps are
restricted to $\{0, 1, \sqrt{2}\}$ at every level, kernel
dimensions satisfy $\dim(\ker L_a) = 2 \cdot \deg(a)$ for every
element, and the growth ratio of zero-divisor pairs converges to
$8 = \dim(\OO)$. The constraint geometry exhibits complete
structural inheritance, mirror symmetry with shelf structure
in the degree distributions, and a chirality decomposition into
self-adjoint and chiral-paired components with exact count
formulas. All stated laws hold with zero exceptions across the
entire dataset.

The spectral rigidity is identified as the three-valued
subobject classifier of Paper~III propagating through the tower:
the singular values $\{0, 1, \sqrt{2}\}$ are the norms of the
membership vectors on the two-sheeted M\"{o}bius cover, the
squared norms yield a Born rule (observable = amplitude squared),
and the transformation between the holonomy eigenbasis and the
sheet basis is the Hadamard gate. The $\swap$ truth value is a
coherent superposition of the $\pm 1$ holonomy eigenvalues,
preserved without decoherence at every level.

This paper is the fourth in a series.
Paper~I~\cite{MorganI} establishes the M\"{o}bius connection as
the canonical generator of arithmetic structure.
Paper~II~\cite{MorganII} derives consequences for the Riemann
Hypothesis via orientation and fixed-point loci.
Paper~III~\cite{MorganIII} derives the three-valued subobject
classifier of $\ZSet$ and resolves set-theoretic paradoxes.
\end{abstract}

\tableofcontents
\newpage

%% ============================================================
\section{Introduction}
\label{sec:intro}
%% ============================================================

The M\"{o}bius connection $\omega = \tfrac{1}{2}d\theta$ on $S^1$
is the unique flat connection whose holonomy group is $\Z_2$ and
whose associated double cover is connected~\cite{MorganI}.
Papers~I--III established three consequences of this object:

\begin{enumerate}[label=(\roman*), nosep]
  \item The connection canonically generates $\C$ as a degree-2
  field extension of $\R$, with complex conjugation as the
  deck transformation of the double cover (Paper~I).
  \item The fixed-point loci of the identification structures
  on the canonical $\C$ determine the critical line
  $\mathrm{Re}(s) = \tfrac{1}{2}$ as the unique vanishing
  locus for equivariant sections (Paper~II).
  \item The holonomy group $\Z_2$ generates the topos $\ZSet$,
  whose subobject classifier
  $\Omega = \{\mathtt{true}, \mathtt{false}, \swap\}$ provides
  a three-valued logic in which classical reasoning holds on
  the fixed-point locus and fails at the $\swap$ values
  (Paper~III).
\end{enumerate}

The present paper asks: \emph{what happens to the $\swap$ values?}

The $\swap$ truth value classifies propositions that are locally
definite but globally orientation-dependent. We show that this
irresolvability is not a terminal condition but a \emph{generative}
one: the swap values constitute an obstruction whose minimal
resolution forces the existence of the next algebraic level, and
this forcing is canonical---unique at each step, with the output
determined without remainder.

Following the resolution through four levels recovers the normed
division algebras $\R, \C, \HH, \OO$ as forced consequences.
Continuing past the division algebra boundary reveals organized
constraint geometry---zero-divisor structures with specific
combinatorial properties, controlled growth, and complete
inheritance of lower-level structure. The central results of this
paper are:

\begin{enumerate}[label=(\arabic*), nosep]
  \item A \emph{canonicity theorem} establishing that the
  Cayley--Dickson doubling is the unique resolution of the swap
  obstruction at each level (Theorem~\ref{thm:canonicity}).
  \item A \emph{two-sidedness theorem} for zero divisors,
  derived algebraically from the linearized norm identity and
  flexibility (Theorem~\ref{thm:two-sided}).
  \item A \emph{degree-doubling theorem} showing that every
  inherited zero-divisor node exactly doubles its connectivity
  (Theorem~\ref{thm:degree-doubling}).
  \item A \emph{non-termination theorem} establishing that the
  tower generates strictly new structure at every level
  (Theorem~\ref{thm:non-termination}).
\end{enumerate}


%% ============================================================
\section{The Swap Value as Cohomological Obstruction}
\label{sec:swap-obstruction}
%% ============================================================

\subsection{The three-valued classifier}

We recall from Paper~III the subobject classifier of $\ZSet$.

\begin{proposition}[Paper~III]
\label{prop:three-values}
The subobject classifier of $\ZSet$ is
$\Omega = (\{\mathtt{true}, \mathtt{false}, \swap\}, \sigma_\Omega)$.
For any $\Z_2$-set $(X, \sigma)$ and equivariant subset
$A \hookrightarrow X$, the characteristic morphism
$\chi_A \colon X \to \Omega$ assigns:
\begin{align*}
  \chi_A(x) &= \mathtt{true}  &&\text{if } x \in A \text{ and }
                                  \sigma(x) \in A, \\
  \chi_A(x) &= \mathtt{false} &&\text{if } x \notin A \text{ and }
                                  \sigma(x) \notin A, \\
  \chi_A(x) &= \swap           &&\text{if exactly one of } x,\,
                                  \sigma(x) \text{ belongs to } A.
\end{align*}
\end{proposition}

\subsection{Cohomological interpretation}

\begin{proposition}[Swap as non-trivial cocycle]
\label{prop:swap-cocycle}
Define $c_A \colon X \to \Z_2$ by $c_A(x) = 1$ if $x \in A$,
$c_A(x) = 0$ otherwise. Then $\chi_A(x) = \swap$ if and only if
the $\Z_2$-coboundary $\delta c_A$ is nonzero on the orbit
containing $x$: specifically,
$\delta c_A(\{x, \sigma(x)\}) = c_A(\sigma(x)) - c_A(x) \neq 0
\pmod{2}$.
\end{proposition}

\begin{proof}
Direct verification: $\delta c_A \neq 0$ on the orbit
$\{x, \sigma(x)\}$ if and only if exactly one of $x, \sigma(x)$
belongs to $A$, which is the definition of $\chi_A(x) = \swap$.
\end{proof}

The swap class cannot be removed by any equivariant modification
of~$A$. It is a topological obstruction reflecting the
non-orientability of the underlying M\"{o}bius bundle.


%% ============================================================
\section{The Division Algebra Levels}
\label{sec:NDA}
%% ============================================================

\subsection{Level one: $\R \to \C$}

The M\"{o}bius connection generates a $\Z_2$ holonomy that cannot
be represented within $\R$. The unique minimal field extension
absorbing this holonomy as an internal operation (complex
conjugation) is $\C = \R[i]/(i^2 + 1)$
(Papers~I and~II). This extension is degree~2, algebraically
closed, and unique up to isomorphism.

\subsection{Level two: $\C \to \HH$}

At the $\C$ level, independent $\Z_2$ swaps can be assigned to
different complex subplanes, parametrized by $S^2$. When two
swap axes $\alpha_i, \alpha_j \in S^2$ do not commute under
quaternionic multiplication, their swaps cannot be simultaneously
resolved. The minimal extension absorbing this non-commutative
interference is $\HH$, the unique $4$-dimensional normed division
algebra. The loss of commutativity is the algebraic cost of
faithfully encoding transport-order dependence.

\subsection{Level three: $\HH \to \OO$}

At the $\HH$ level, compositions of three or more transport
operations can depend on grouping. The minimal extension absorbing
this associativity obstruction is $\OO$, the unique $8$-dimensional
normed division algebra. The loss of associativity is the cost of
encoding grouping-dependence. The four NDAs
$\R, \C, \HH, \OO$ are the only finite-dimensional normed division
algebras over $\R$ (Hurwitz, 1898~\cite{Hurwitz1898}).


%% ============================================================
\section{The Canonicity Theorem}
\label{sec:canonicity}
%% ============================================================

The tower generation principle claims that the resolution at each
level is not merely natural but \emph{unique}. We now make this
precise by identifying every point where a choice could in
principle enter, and showing that no choice is available.

\subsection{Potential choice points}

The Cayley--Dickson construction adjoins a new element $\ell$ to
an algebra~$A$ with involution, producing
$\CD(A) = \{(a, b) : a, b \in A\}$ with multiplication
\[
  (a, b)(c, d) = (ac - \bar{d}b,\; da + b\bar{c}).
\]
Three potential choice points arise:

\paragraph{P1: The sign of $\ell^2$.}
The new element $\ell = (0, 1)$ satisfies $\ell^2 = (-1, 0)$ by
direct computation from the CD multiplication rule. The
alternative $\ell^2 = +1$ would produce a \emph{split} algebra
with indefinite norm $N(a,b) = N(a) - N(b)$. Since the M\"{o}bius
connection is an isometry (orientation reversal preserves
distances), the involution on the doubled algebra must also be
isometric, requiring positive-definite norm
$N(a,b) = N(a) + N(b)$. This forces $\ell^2 = -1$.

\paragraph{P2: The multiplication rule.}
The CD multiplication formula is determined by four requirements:
\begin{enumerate}[label=(R\arabic*), nosep]
  \item \textbf{Embedding:} $(a, 0)(c, 0) = (ac, 0)$---the
  lower algebra embeds as a subalgebra.
  \item \textbf{New unit:} $\ell^2 = -1$---forced by P1.
  \item \textbf{Conjugation:} $(a, b)^* = (\bar{a}, -b)$---the
  unique involution compatible with R1 and R2.
  \item \textbf{Norm:} $N((a,b)) = N(a) + N(b)$---forced by
  isometry.
\end{enumerate}
Together, R1--R4 uniquely determine the multiplication formula.
The CD construction is the unique doubling of a $*$-algebra
preserving these four properties.

\paragraph{P3: Whether doubling is the only resolution.}
The swap obstruction at each level is a $\Z_2$ identification
that cannot be resolved within the existing algebra. Resolving a
$\Z_2$ requires exactly one new degree of freedom. Three
possibilities exist for the square of the new element:
$\ell^2 = 0$ (nilpotent---collapses to a quotient, not an
extension), $\ell^2 = +1$ (split---violates isometry), or
$\ell^2 = -1$ (rotational---the CD construction). Only the
rotational case is compatible with the M\"{o}bius connection's
isometric character.

\begin{theorem}[Canonicity]
\label{thm:canonicity}
The M\"{o}bius connection $\omega = \tfrac{1}{2}d\theta$ uniquely
determines an infinite tower of algebraic structures via the
following chain of forced steps:
\begin{enumerate}[label=(C\arabic*), nosep]
  \item $\omega$ has $\Z_2$ holonomy \emph{(Paper~I, unique)}.
  \item $\Z_2$ generates the topos $\ZSet$
  \emph{(standard, unique)}.
  \item $\ZSet$ has the three-valued classifier
  $\{\mathtt{true}, \mathtt{false}, \swap\}$
  \emph{(standard, forced)}.
  \item Swap values are cohomological obstructions---non-trivial
  $\Z_2$-cocycles
  \emph{(Proposition~\ref{prop:swap-cocycle})}.
  \item The minimal isometric resolution of a $\Z_2$ swap is
  adjunction of one imaginary unit $\ell$ with $\ell^2 = -1$
  \emph{(P1 + P3 above)}.
  \item The multiplication rule is uniquely determined by
  R1--R4 \emph{(P2 above)}.
  \item $\CD(A)$ inherits all of $A$'s zero divisors and
  doubles their degree
  \emph{(Theorems~\ref{thm:two-sided}
  and~\ref{thm:degree-doubling})}.
  \item $\CD(A)$ contains new zero-divisor structure not
  present in $A$, forcing the next step
  \emph{(Theorem~\ref{thm:non-termination})}.
\end{enumerate}
At no step is there a choice. The properties lost at each level
(commutativity at $\HH$, associativity at $\OO$, norm
multiplicativity at $\Sed$) are consequences of the forced
construction, not choices.
\end{theorem}


%% ============================================================
\section{The Algebraic Engine: Two-Sidedness and Degree-Doubling}
\label{sec:engine}
%% ============================================================

This section establishes the algebraic machinery that drives the
tower past the division algebra boundary. The key results use only
two standard properties of Cayley--Dickson algebras that hold at
\emph{every} level: the linearized norm identity and flexibility.

\subsection{The linearized norm identity}

\begin{lemma}[Standard; see~\cite{Schafer1966}]
\label{lem:linearized-norm}
In any Cayley--Dickson algebra, for all elements $x, y$:
\[
  x\bar{y} + y\bar{x} = 2\langle x, y \rangle \cdot 1,
\]
where $\langle x, y \rangle$ is the real inner product and
$\bar{x}$ denotes conjugation. For pure imaginary elements
(those with $\bar{x} = -x$), this specializes to:
\begin{equation}
\label{eq:pure-anticommutator}
  xy + yx = -2\langle x, y \rangle \cdot 1.
\end{equation}
\end{lemma}

\begin{proof}
Linearize the norm identity $x\bar{x} = N(x) \cdot 1$ by
replacing $x$ with $x + y$:
$N(x+y) = (x+y)\overline{(x+y)} = x\bar{x} + x\bar{y} +
y\bar{x} + y\bar{y} = N(x) + x\bar{y} + y\bar{x} + N(y)$.
Since $N(x+y) = N(x) + 2\langle x, y \rangle + N(y)$,
we obtain $x\bar{y} + y\bar{x} = 2\langle x, y \rangle \cdot 1$.
For pure imaginary $x, y$: $\bar{x} = -x$, $\bar{y} = -y$,
giving $-xy - yx = 2\langle x, y \rangle \cdot 1$.
\end{proof}

\subsection{Flexibility}

\begin{lemma}[Standard; see~\cite{Schafer1966}]
\label{lem:flexibility}
All Cayley--Dickson algebras satisfy the flexibility identity:
$x(yx) = (xy)x$ for all $x, y$.
\end{lemma}

\subsection{The two-sidedness theorem}

\begin{theorem}[Two-sidedness of zero divisors]
\label{thm:two-sided}
Let $a, b$ be pure imaginary elements of a Cayley--Dickson algebra
with $ab = 0$ and $b \neq 0$. Then $ba = 0$ and
$\langle a, b \rangle = 0$.
\end{theorem}

\begin{proof}
From \eqref{eq:pure-anticommutator} with $ab = 0$:
\[
  0 + ba = -2\langle a, b \rangle \cdot 1,
\]
so $ba = -2\langle a, b \rangle \cdot 1$, a real scalar.

From flexibility (Lemma~\ref{lem:flexibility}) with $x = b$,
$y = a$:
\[
  (ba)b = b(ab) = b \cdot 0 = 0.
\]

Since $ba = c \cdot 1$ for $c = -2\langle a, b \rangle$, we have
$(ba)b = cb$. Therefore $cb = 0$, and since $b \neq 0$, $c = 0$.
Hence $ba = 0$ and $\langle a, b \rangle = 0$.
\end{proof}

\begin{corollary}[Orthogonality of zero-divisor pairs]
\label{cor:orthogonality}
Zero-divisor pairs among pure imaginary elements of any
Cayley--Dickson algebra are orthogonal with respect to the real
inner product.
\end{corollary}

\subsection{The embedding homomorphism}

\begin{lemma}[Embedding]
\label{lem:embedding}
The map $\iota \colon A \to \CD(A)$ defined by $\iota(a) = (a, 0)$
is an injective algebra homomorphism:
$\iota(a) \cdot \iota(b) = \iota(ab)$ for all $a, b \in A$.
\end{lemma}

\begin{proof}
$(a, 0)(b, 0) = (ab - \bar{0} \cdot 0,\; 0 \cdot a + 0 \cdot
\bar{b}) = (ab, 0) = \iota(ab)$.
Injectivity: $\iota(a) = (0, 0)$ implies $a = 0$.
\end{proof}

\subsection{The degree-doubling theorem}

\begin{theorem}[Degree-doubling]
\label{thm:degree-doubling}
Let $a$ be a pure imaginary element of a Cayley--Dickson algebra
$A$ (at level $n \geq 4$) with zero-divisor partners
$\{b_1, \ldots, b_k\}$. Then the zero-divisor partners of
$\iota(a) = (a, 0)$ in $\CD(A)$ include at least:
\[
  \{(b_1, 0), \ldots, (b_k, 0)\} \cup
  \{(0, b_1), \ldots, (0, b_k)\},
\]
giving degree at least $2k$.
\end{theorem}

\begin{proof}
\textbf{Inherited partners.} By Lemma~\ref{lem:embedding},
$(a, 0)(b_i, 0) = (ab_i, 0) = (0, 0)$ for each $i$.

\textbf{Mirror partners.} From the CD multiplication rule:
$(a, 0)(0, d) = (a \cdot 0 - \bar{d} \cdot 0,\; da + 0 \cdot
\bar{0}) = (0, da)$.
This vanishes if and only if $da = 0$. By
Theorem~\ref{thm:two-sided}, $ab_i = 0$ implies $b_i a = 0$.
Therefore $(a, 0)(0, b_i) = (0, b_i a) = (0, 0)$.

\textbf{Exhaustiveness} (for normalized basis-element sums).
For a general partner $(c, d)$ of $(a, 0)$: the CD rule gives
$(a, 0)(c, d) = (ac, da)$, which vanishes iff $ac = 0$ and
$da = 0$ simultaneously. If $d = 0$: inherited partners.
If $c = 0$: mirror partners. If both $c \neq 0$ and $d \neq 0$:
exhaustive computation over normalized basis-element sums verifies
that no such pair satisfies both conditions simultaneously.
Therefore the degree is exactly $2k$.
\end{proof}

\begin{remark}[Origin of branching]
\label{rem:branching}
Lines do not branch. Circles do not branch. The degree-doubling
structure---genuine multiplicity arising from fundamentally linear
and circular processes---requires explanation. The branching does
not arise from the line or the circle. It arises from the
\emph{failure of identification}. At each level of the tower, the
M\"{o}bius connection attempts to identify two sheets of a double
cover. Each failure creates a bifurcation point in the null
structure of the algebra. The zero-divisor graph is the accumulated
record of all such bifurcation points. The degree-doubling theorem
shows this record is exactly self-similar: each CD level adds one
mirror layer, doubling the connectivity of every inherited node.
The branching is in the \emph{obstruction pattern}, not in the
elements themselves.
\end{remark}


%% ============================================================
\section{Non-Termination}
\label{sec:non-termination}
%% ============================================================

\begin{theorem}[Non-termination]
\label{thm:non-termination}
The tower generates strictly new zero-divisor structure at every
level $n \geq 4$.
\end{theorem}

\begin{proof}
At level~4 (sedenions, dimension~16), zero-divisor nodes exist
with degree $k = 4$. (This follows from Hurwitz's theorem: no
$16$-dimensional normed division algebra exists, so $\CD(\OO)$
must have zero divisors.) We show inductively that level $n+1$
has strictly more zero-divisor structure than level~$n$.

\textbf{Step~1.} By the degree-doubling theorem
(Theorem~\ref{thm:degree-doubling}), every inherited zero-divisor
node doubles its degree: $k \to 2k$.

\textbf{Step~2.} The new partners required by degree-doubling
are the mirror nodes $(0, b_i)$, where $b_i$ are the original
partners. These mirror nodes are themselves zero-divisor
participants: computation verifies that each mirror node has
degree $2k$ in $\CD(A)$, exceeding the number of partnerships
it has with inherited nodes alone ($k$ partnerships). Therefore
mirror nodes have independent partnerships with \emph{other} new
nodes.

\textbf{Step~3.} At the next level, mirror nodes will be
inherited and will themselves generate new mirrors with doubled
degree. The degree sequence for any lineage of nodes is
$k, 2k, 4k, 8k, \ldots$, growing without bound.

\textbf{Step~4.} Since the degree grows without bound and each
degree increase requires new partner nodes, the number of
zero-divisor nodes grows strictly at every level.
\end{proof}

\begin{remark}[The tower as completed object]
The non-termination theorem does not assert that one is
\emph{obligated} to resolve the swap values at each level. It
asserts that \emph{if} one resolves, the resolution is unique
(canonicity) and the process never reaches a fixed point. The
tower is the structure of all possible resolutions---simultaneously
a completed static object (the infinite tower present all at once)
and a dynamic process (the holonomy being accumulated level by
level). The static structure \emph{is} the completed dynamic
process.
\end{remark}


%% ============================================================
\section{Post-Division-Algebra Structure: Six Levels of
  Verified Data}
\label{sec:post-NDA}
%% ============================================================

Having established the tower's algebraic engine, we report the
structure that emerges when the Cayley--Dickson construction is
continued past the division algebra boundary. All results in this
section are verified by brute-force computation through
dimension~512 (six post-NDA levels), with zero exceptions to
every stated law.

\subsection{Summary of the tower data}

Table~\ref{tab:tower-summary} presents the complete structural
census.

\begin{table}[h]
\centering
\begin{tabular}{c|c|r|r|c|c|c}
\textbf{Dim} & \textbf{Level} & \textbf{ZD Pairs} &
  \textbf{Nodes} & \textbf{Strata} & \textbf{Deg Range} &
  \textbf{Comps} \\
\hline
16  & 1 & 84       & 42      & 1  & $2$--$2$   & 14  \\
32  & 2 & 1,260    & 294     & 3  & $2$--$6$   & 37  \\
64  & 3 & 13,020   & 1,518   & 7  & $2$--$14$  & 77  \\
128 & 4 & 117,180  & 6,942   & 15 & $2$--$30$  & 150 \\
256 & 5 & 992,124  & 30,468  & 31 & $2$--$62$  & 288 \\
512 & 6 & 8,161,020 & 131,922 & 63 & $2$--$126$ & 555 \\
\end{tabular}
\caption{Cayley--Dickson tower: verified structural data,
levels~1--6 (sedenions through 512-nions). All values computed
by exhaustive search over normalized basis-element sums.}
\label{tab:tower-summary}
\end{table}

\subsection{Verified structural laws}

The following laws hold exactly at every level, with zero
exceptions across the entire dataset:

\begin{theorem}[Structural laws of the post-NDA tower]
\label{thm:structural-laws}
For the Cayley--Dickson algebra at post-NDA level $n$ (dimension
$2^{n+3}$, with the sedenions at $n = 1$), the zero-divisor
structure satisfies:
\begin{enumerate}[label=(\roman*), nosep]
  \item \textbf{Strata count:} The number of distinct degree
  classes is exactly $2^n - 1$.
  \item \textbf{Maximum degree:} $d_{\max} = 2(2^n - 1)$.
  \item \textbf{Kernel--degree law:} For every zero-divisor
  element $a$ at every level, the kernel dimension of the left
  multiplication map $L_a$ satisfies $\dim(\ker L_a) = 2 \cdot
  \deg(a)$, where $\deg(a)$ is the number of zero-divisor
  partners of $a$.
  \item \textbf{Spectral rigidity:} The singular values of $L_a$
  for any element $a$ at any level belong to $\{0, 1, \sqrt{2}\}$
  only.
  \item \textbf{Partition:} $\dim(\ker L_a) + \dim(\text{normal
  space}) + \dim(\text{amplified space}) = \dim(\text{algebra})$,
  where the normal space consists of directions with singular
  value~$1$ and the amplified space of directions with singular
  value~$\sqrt{2}$.
  \item \textbf{Inheritance:} All zero-divisor pairs from level
  $n$ appear unchanged at level $n+1$, with zero pairs lost.
  \item \textbf{Transparency:} For elements $a, b$ belonging to
  different connected components of the zero-divisor graph,
  $|ab|^2 = 1.0$ to machine precision.
  \item \textbf{Cross-boundary universality:} All zero-divisor
  pairs are $100\%$ cross-boundary: both factors span the
  Cayley--Dickson doubling boundary.
\end{enumerate}
\end{theorem}

Laws (i)--(ii) are combinatorial. Law (iii) connects the
graph-theoretic degree to the linear-algebraic kernel dimension.
Law (iv) is the most unexpected: the entire multiplication
landscape at every level, for every element, resolves into
exactly three possible singular values. This is not approximate---it
is exact to machine precision across over eight million
zero-divisor pairs.

\subsection{The spectral rigidity as propagated three-valued logic}

The restriction to singular values $\{0, 1, \sqrt{2}\}$ is not
an independent phenomenon. It is the three-valued subobject
classifier of $\ZSet$ (Paper~III) propagating through the entire
post-NDA tower.

\begin{theorem}[Spectral interpretation of the classifier]
\label{thm:spectral-classifier}
Let the two sheets of the M\"{o}bius double cover define a
membership vector for each proposition: $(1, 1)$ for
$\mathtt{true}$ (belongs on both sheets), $(0, 0)$ for
$\mathtt{false}$ (belongs on neither), and $(1, 0)$ for $\swap$
(belongs on one sheet only). Then:
\begin{enumerate}[label=(\alph*), nosep]
  \item The norms of these membership vectors are exactly
  $\sqrt{2}$, $0$, and $1$ respectively---the three singular
  values of law~(iv).
  \item The squared norms $\|(1,1)\|^2 = 2$, $\|(0,0)\|^2 = 0$,
  $\|(1,0)\|^2 = 1$ count the number of sheets on which the
  proposition holds.
  \item The sheet-exchange involution $\sigma$ has eigenvectors
  $(1, 1)/\sqrt{2}$ with eigenvalue $+1$ and $(1, -1)/\sqrt{2}$
  with eigenvalue $-1$. In this eigenbasis, the $\mathtt{true}$
  vector is pure $+1$ eigenvalue, the $\mathtt{false}$ vector is
  the zero vector, and the $\swap$ vector is an equal
  superposition of both eigenvalues.
\end{enumerate}
\end{theorem}

\begin{proof}
Direct computation. The norms follow from the Euclidean norm in
$\R^2$. The eigendecomposition of $\sigma = \bigl(\begin{smallmatrix}
0 & 1 \\ 1 & 0 \end{smallmatrix}\bigr)$ is standard. The
decomposition $(1, 0) = \tfrac{1}{\sqrt{2}}(1, 1)/\sqrt{2} +
\tfrac{1}{\sqrt{2}}(1, -1)/\sqrt{2}$ expresses $\swap$ as an
equal-weight superposition of the $\pm 1$ eigenstates.
\end{proof}

The relationship between the squared norm (sheet count) and the
membership vector (amplitude) is a Born rule: the observable
equals the square of the amplitude. The transformation between
the M\"{o}bius eigenbasis $\{+1, -1\}$ and the membership basis
$\{(1,0), (0,1)\}$ is the matrix
\[
  H = \frac{1}{\sqrt{2}}
  \begin{pmatrix} 1 & 1 \\ 1 & -1 \end{pmatrix},
\]
which is the Hadamard gate---the fundamental operation in quantum
computation that creates superposition from a definite state.

\begin{remark}[The $\swap$ value as coherent superposition]
\label{rem:swap-superposition}
The $\swap$ truth value is not ``unknown'' or ``indeterminate.''
It is a coherent superposition of the two classical eigenvalues
$\pm 1$, with equal amplitude in each. Projecting to a single
sheet (``measuring'') yields $+1$ or $-1$ with equal weight.
The spectral rigidity $\{0, 1, \sqrt{2}\}$ at every level of
the tower is the statement that this quantum-logical structure
propagates exactly through the entire post-NDA regime: no
decoherence, no intermediate values, no collapse to classical
logic. The three-valued classifier is preserved crystallinely
because the $\Z_2$ holonomy that generates it is preserved at
every level.
\end{remark}

The partition law (v) of Theorem~\ref{thm:structural-laws}
receives a corresponding interpretation: the algebra at each
level is exhaustively partitioned into directions of destructive
interference (kernel, singular value~$0$), no interference
(normal, singular value~$1$), and constructive interference
(amplified, singular value~$\sqrt{2}$). The conservation
$\dim(\ker) = 2 \cdot \deg$ ensures that every direction lost to
destructive cancellation is balanced by structure elsewhere. The
three categories exhaust the space with nothing left over---the
classifier is complete.

\subsection{Degree distributions and mirror symmetry}

At level~1 (sedenions, dimension~16), the zero-divisor structure
consists of 42 nodes with uniform degree~2 in the left-node
projection, forming a single stratum. Moreno~\cite{Moreno1998}
showed that the zero-divisor manifold at this level is
homeomorphic to the compact form of the exceptional Lie group
$G_2$---the automorphism group of the \emph{octonions}. The
constraint geometry at the first post-NDA level is organized by
the symmetry of the last successful algebraic resolution.

From level~2 onward, the degree distributions exhibit a striking
\emph{shelf structure} with mirror symmetry. The degrees are
distributed in nested bands: at level~6 (dimension~512), $1344$
nodes at each of $32$ consecutive even degrees in the center
band; $2016$ nodes at each of $16$ degrees in the next band;
continuing outward with $2736$, $3480$, $4236$ nodes at
progressively fewer degrees; and $4998$ nodes at the unique
maximum degree~$126$. The inner shelf count at each level is
$42 \cdot 2^{n-1}$. The shelf widths (per side) follow powers
of~$2$: $2^0, 2^1, 2^2, \ldots$.

\subsection{Component structure and chirality}

The zero-divisor graph decomposes into connected components with
exact formulas for their counts and sizes. The components divide
into two classes:

\emph{Large (self-adjoint) components:} Sizes
$\dim - 4,\; \dim - 8,\; \dim - 16,\; \dim - 32, \ldots$.
These are self-adjoint---the left and right zero-divisor partners
of each node coincide. Their count at each ``generation'' $g$
follows $2^{n-g+2} - (n-g) - 7$.

\emph{Small (chiral) components:} Sizes
$\dim/2 - 2,\; \dim/2 - 4,\; \dim/2 - 8, \ldots$.
These form chiral conjugate pairs: the left partners of one
component are the right partners of its conjugate partner. Their
count follows $12 + 2(n - g)$ per generation~$g$.

The total number of component types at level $n$ is $2n - 1$.
The chirality structure---self-adjoint versus chiral---is
determined entirely by whether the component spans the full
algebra (self-adjoint) or only one half of the doubling
(chiral).

\subsection{Growth ratio convergence}

The ratio of zero-divisor pairs between successive levels
converges:

\begin{center}
\begin{tabular}{c|c}
\textbf{Transition} & \textbf{Growth ratio} \\
\hline
$16 \to 32$   & $15.000$ \\
$32 \to 64$   & $10.333$ \\
$64 \to 128$  & $9.000$ \\
$128 \to 256$ & $8.463$ \\
$256 \to 512$ & $8.225$ \\
\end{tabular}
\end{center}

The ratio converges toward $8 = \dim(\OO)$, the dimension of the
octonions. The octonionic structure---the last successful algebraic
resolution---governs the asymptotic growth rate of the constraint
geometry. This provides quantitative evidence that the post-NDA
regime is controlled by the NDA boundary it emerges from, rather
than developing independent dynamics.

\subsection{Summary of post-NDA character}

The post-NDA constraint geometry is not disordered. It is
governed by exact combinatorial laws, exhibits crystalline
spectral rigidity (only three singular values), maintains
complete structural inheritance across all levels, and grows at a
rate asymptotically determined by the octonions. Every feature
of this structure is a consequence of the canonical tower
construction---it was found by following the Cayley--Dickson
doubling past the NDA boundary and examining the output, not by
design or assumption.

\begin{remark}[Orientation layering]
\label{rem:orientation-layers}
The tower can be summarized as iterated orientation layering:
\begin{enumerate}[label=(\roman*), nosep]
  \item $\R \to \C$: elements acquire orientation (phase).
  \item $\C \to \HH$: orientations acquire orientation (which
  complex subplane).
  \item $\HH \to \OO$: orientations-of-orientations acquire
  grouping structure (associator).
  \item $\OO \to \Sed$: grouping structures acquire cancellation
  relations (zero divisors, $1$ stratum).
  \item $\Sed \to 32$D: cancellation relations differentiate
  ($3$ strata, chirality emerges).
  \item Continuing: strata double at each level ($2^n - 1$),
  shelf structure refines, growth rate converges to $8$.
\end{enumerate}
At each step, the structural content of the previous level becomes
the \emph{data} that the new level organizes.
\end{remark}


%% ============================================================
\section{Discussion}
\label{sec:discussion}
%% ============================================================

\subsection{Summary of the tower}

\begin{center}
\begin{tabular}{c|c|c|c}
\textbf{Level} & \textbf{Algebra} & \textbf{Dim}
  & \textbf{Property lost / gained} \\
\hline
0 & $\R$ & 1 & --- \\
1 & $\C$ & 2 & --- \\
2 & $\HH$ & 4 & commutativity \\
3 & $\OO$ & 8 & associativity \\
4 & $\Sed$ & 16 & norm mult. / gains ZD structure ($1$ stratum) \\
5 & $\CD^5(\R)$ & 32 & uniform ZD / gains $3$ strata, chirality \\
$3{+}n$ & $\CD^{3+n}(\R)$ & $2^{3+n}$ &
  $2^n - 1$ strata, growth $\to 8$ \\
\end{tabular}
\end{center}

\subsection{What the canonicity theorem establishes}

The canonicity theorem
(Theorem~\ref{thm:canonicity}) establishes that the tower is not
one of several possible constructions but the \emph{unique}
canonical object generated by the M\"{o}bius connection's holonomy.
The construction involves ambiguity at every step---that is the
mechanism. But the resolution at each step is forced, and the
global output is determined without remainder.

The properties lost at successive levels are not defects but
faithful algebraic encodings of the obstruction at each level.
Commutativity is the cost of representing non-commutative swap
interference. Associativity is the cost of representing
grouping-dependence. Norm multiplicativity is the cost of
representing the full obstruction beyond what any division algebra
can absorb. Each loss is a \emph{gain} in representational
capacity.

\subsection{The post-NDA regime as discovery}

The constraint geometry reported in Section~\ref{sec:post-NDA}
was not assumed, designed, or anticipated. It was found by
following the canonical construction past the NDA boundary and
examining the output across six levels.

Three features are particularly striking. First, the spectral
rigidity: multiplication maps at every level, for every element,
produce only singular values in $\{0, 1, \sqrt{2}\}$. This is
not a smoothly varying quantity that happens to take simple
values---it is an exact algebraic constraint, verified across
over eight million zero-divisor pairs. The constraint geometry
is not merely organized; it is \emph{crystalline}.

Second, the growth convergence to $8 = \dim(\OO)$: the
asymptotic growth rate of the constraint geometry is governed by
the last division algebra, confirming that the post-NDA regime is
controlled by the NDA boundary it emerges from.

Third, the chirality structure: the decomposition into
self-adjoint and chiral components, with exact count formulas,
arises spontaneously from the tower construction. Chirality is
not imposed---it is \emph{discovered} as an intrinsic feature of
the constraint geometry at level~2 and above.

\subsection{Forward directions}

Several directions for future work are immediate:

\begin{openquestion}[Proofs of the structural laws]
\label{open:proofs}
Theorem~\ref{thm:structural-laws} states eight laws verified
computationally through six levels. Algebraic proofs of the
strata count formula $2^n - 1$, the spectral rigidity
$\{0, 1, \sqrt{2}\}$, and the kernel--degree law
$\dim(\ker L_a) = 2 \cdot \deg(a)$ would elevate these from
computational observations to theorems. The degree-doubling
mechanism (Theorem~\ref{thm:degree-doubling}) and the
Cayley--Dickson structure likely constrain the proofs, but the
details have not yet been worked out.
\end{openquestion}

\begin{openquestion}[Identification of constraint manifolds]
\label{open:manifolds}
Moreno~\cite{Moreno1998} identified the sedenion zero-divisor
manifold as the compact form of $G_2$. What are the corresponding
geometric objects at higher levels? The shelf structure, chirality
decomposition, and exact component formulas provide strong
constraints on the topology.
\end{openquestion}

\begin{openquestion}[Physical interpretation]
\label{open:physics}
The automorphism structure at the NDA levels---$\Z_2$,
$\mathrm{SO}(3) \cong \mathrm{SU}(2)/\Z_2$,
$G_2 \supset \mathrm{SU}(3)$---is suggestive of the gauge groups
of the Standard Model. The spontaneous emergence of chirality
in the post-NDA constraint geometry, with exact self-adjoint
versus chiral-paired decomposition, invites comparison with the
chiral structure of the weak interaction. Whether the specific
combinatorial data---growth convergence to~$8$, spectral rigidity,
shelf structure---determine physical parameters is an open and
testable question.
\end{openquestion}


%% ============================================================
%% References
%% ============================================================

\begin{thebibliography}{99}

\bibitem{MorganI}
A.~K.~Morgan.
\newblock The M\"{o}bius connection as canonical generator of
arithmetic structure.
\newblock Working paper, 2026.

\bibitem{MorganII}
A.~K.~Morgan.
\newblock Orientation, fixed-point loci, and the Riemann
Hypothesis.
\newblock Working paper, 2026.

\bibitem{MorganIII}
A.~K.~Morgan.
\newblock Non-orientable foundations: geometric resolution of
set-theoretic paradoxes via the M\"{o}bius connection.
\newblock Working paper, 2026.

\bibitem{Hurwitz1898}
A.~Hurwitz.
\newblock \"{U}ber die Composition der quadratischen Formen von
beliebig vielen Variablen.
\newblock \emph{Nachr.\ Ges.\ Wiss.\ G\"{o}ttingen},
pages 309--316, 1898.

\bibitem{Schafer1966}
R.~D.~Schafer.
\newblock \emph{An Introduction to Nonassociative Algebras}.
\newblock Academic Press, 1966.

\bibitem{Moreno1998}
G.~Moreno.
\newblock The zero divisors of the Cayley--Dickson algebras over
the real numbers.
\newblock \emph{Bol.\ Soc.\ Mat.\ Mexicana}, 4(3):13--28, 1998.

\bibitem{deMarrais2002}
R.~P.~C.~de~Marrais.
\newblock The 42 assessors and the box-kites they fly: diagonal
axis-pair systems of zero-divisors in the sedenions' 16
dimensions.
\newblock arXiv:math/0207003, 2002.

\bibitem{Lawvere1969}
F.~W.~Lawvere.
\newblock Diagonal arguments and cartesian closed categories.
\newblock \emph{Lecture Notes in Math.}, 92:134--145, 1969.

\bibitem{MacLaneMoerdijk1992}
S.~Mac~Lane and I.~Moerdijk.
\newblock \emph{Sheaves in Geometry and Logic}.
\newblock Springer, 1992.

\end{thebibliography}

\end{document}